\begin{helper}
Let \(V\) be finite and let \(\nu\) be a measure on activation--priority
pairs \((A,\pi)\in \mathcal P(V)\times \mathbb R^V\).  Fix \(0\le p\le1\),
distinct vertices \(u,v\in V\), and disjoint finite sets
\(U,T\subseteq V\setminus\{u,v\}\).  Assume the activation atom law
\[
\nu(A=A_0)=p^{|A_0|}(1-p)^{|V|-|A_0|}
\]
for every \(A_0\subseteq V\).  Assume also that, on every activation atom,
the activated priorities are supported in the unit cube, and that every lower
rectangle satisfies
\[
\nu(A=A_0,\ 0\le \pi(q)\le t_q\text{ for all }q\in V)
=\nu(A=A_0)\prod_{q\in V}t_q
\]
whenever \(0\le t_q\le1\) for all \(q\).

Then the ordered two-coordinate chamber
\[
\begin{aligned}
E_{u,v}(U,T)=\{(A,\pi):\;&u,v\in A,\ \pi(v)\le \pi(u),\
\pi(q)\in[0,1]\text{ for every }q\in A,\\
& z\in U\cap A\Rightarrow \pi(z)<\pi(u),\
z\in T\cap A\Rightarrow \pi(z)<\pi(v)\}
\end{aligned}
\]
has mass
\[
\nu(E_{u,v}(U,T))
=\operatorname{ofReal}\!\left(
\int_0^1\int_0^a
p^2\bigl((1-p)+pa\bigr)^{|U|}
\bigl((1-p)+pb\bigr)^{|T|}\,db\,da\right).
\]
Moreover the diagonal chamber
\[
\{u,v\in A,\ \pi(q)\in[0,1]\text{ for }q\in A,\ \pi(u)=\pi(v)\}
\]
has \(\nu\)-measure \(0\), and for every \(z\in U\cup T\) the boundary event
\[
\{u,v,z\in A,\ \pi(q)\in[0,1]\text{ for }q\in A,\ 
\pi(z)=\pi(u)\text{ or }\pi(z)=\pi(v)\}
\]
has \(\nu\)-measure \(0\).
\end{helper}
\begin{proof}
Write \(S=U\cup T\) and \(C=S\cup\{u,v\}\).  The assumptions
\(u,v\notin S\), \(u\ne v\), and \(U\cap T=\emptyset\) imply that the four
sets \(U,T,\{u\},\{v\}\) are pairwise disjoint.  We first record the exact
finite product consequence of the hypotheses on a fixed activation atom.

Fix \(A_0\subseteq V\), and let \(\alpha(A_0)=
p^{|A_0|}(1-p)^{|V|-|A_0|}\).  By the activation atom law,
\(\nu\{A=A_0\}=\operatorname{ofReal}(\alpha(A_0))\).  Since \(0\le p\le1\),
\(\alpha(A_0)\ge0\).  Apply the lower-rectangle hypothesis with the constant
threshold \(1\).  It gives
\[
\nu\{A=A_0,\ 0\le \pi(q)\le1\hbox{ for all }q\in V\}
=\nu\{A=A_0\}.
\]
Thus the part of the atom outside the full cube \([0,1]^V\) has
\(\nu\)-measure \(0\).  On the cube, the same lower-rectangle hypothesis says
that the atom-restricted measure agrees on every anchored rectangle
\(\prod_q[0,t_q]\), \(0\le t_q\le1\), with
\(\operatorname{ofReal}(\alpha(A_0))\) times finite product Lebesgue measure.
The anchored rectangles are closed under finite intersections and generate
the Borel sigma-algebra of \([0,1]^V\).  Both finite measures have the same
total mass \(\operatorname{ofReal}(\alpha(A_0))\).  The finite
pi-lambda uniqueness theorem therefore gives, for every Borel set
\(B\subseteq[0,1]^V\),
\[
\nu\{A=A_0,\ \pi\in B\}
=\operatorname{ofReal}(\alpha(A_0)\lambda_V(B)),
\tag{1}
\]
where \(\lambda_V\) is product Lebesgue measure on the cube.  In particular,
the fixed half-open coordinate cells used in this uniqueness step have the
same masses as in \noderef{SamplingFiniteCoordinateMarginalCellLaw}, and that
node is the registered finite-cell form of the lower-rectangle calculation
from \noderef{SamplingFinitePriorityCutProductFormula}.

Now partition the ordered chamber by the activation pattern on \(S\) and by
the arbitrary activation choices outside \(C\).  For \(B_U\subseteq U\),
\(B_T\subseteq T\), and \(R\subseteq V\setminus C\), put
\[
A_{B_U,B_T,R}=\{u,v\}\cup B_U\cup B_T\cup R .
\]
The corresponding chamber piece is
\[
\begin{aligned}
F(B_U,B_T,R)=\{A=A_{B_U,B_T,R},\;&0\le\pi(q)\le1\hbox{ for all }q,\\
&\pi(v)\le\pi(u),\
z\in B_U\Rightarrow \pi(z)<\pi(u),\
z\in B_T\Rightarrow \pi(z)<\pi(v)\}.
\end{aligned}
\]
The pieces \(F(B_U,B_T,R)\) are pairwise disjoint, and their finite union
differs from \(E_{u,v}(U,T)\) only by atomwise cube complements of measure
zero.  Hence the measure of \(E_{u,v}(U,T)\) is the finite sum of the measures
of these pieces.

For a fixed triple \((B_U,B_T,R)\), formula (1) reduces the piece mass to the
Lebesgue volume in the priority coordinates.  Set \(a=\pi(u)\) and
\(b=\pi(v)\).  By Tonelli's theorem for the finite product cube,
\[
\lambda_V(F(B_U,B_T,R)\hbox{ inside the atom})
=\int_0^1\int_0^a a^{|B_U|}b^{|B_T|}\,db\,da .
\]
Indeed, after \(a,b\) are fixed with \(0\le b\le a\le1\), each coordinate in
\(B_U\) may range independently over an interval of length \(a\), each
coordinate in \(B_T\) over an interval of length \(b\), and every remaining
coordinate has interval length \(1\).  Therefore
\[
\nu(F(B_U,B_T,R))
=\operatorname{ofReal}\!\left(
\alpha(A_{B_U,B_T,R})
\int_0^1\int_0^a a^{|B_U|}b^{|B_T|}\,db\,da\right).
\tag{2}
\]

Summing (2) first over \(R\subseteq V\setminus C\) removes the outside
coordinates:
\[
\sum_R p^{|R|}(1-p)^{|V\setminus C|-|R|}
=(p+(1-p))^{|V\setminus C|}=1 .
\]
Thus the remaining real density is
\[
p^2
\sum_{B_U\subseteq U}\sum_{B_T\subseteq T}
p^{|B_U|+|B_T|}(1-p)^{|S|-|B_U|-|B_T|}
\int_0^1\int_0^a a^{|B_U|}b^{|B_T|}\,db\,da .
\]
All sums are finite and the integrands are nonnegative, so finite additivity
and Tonelli allow us to interchange the finite sums with the two integrals.
The \(U\)-sum factors as
\[
\sum_{B_U\subseteq U}p^{|B_U|}(1-p)^{|U|-|B_U|}a^{|B_U|}
=((1-p)+pa)^{|U|},
\]
and the \(T\)-sum similarly gives \(((1-p)+pb)^{|T|}\).  Hence
\[
\nu(E_{u,v}(U,T))
=\operatorname{ofReal}\!\left(
\int_0^1\int_0^a
p^2((1-p)+pa)^{|U|}((1-p)+pb)^{|T|}\,db\,da\right),
\]
which is the first claimed identity.

It remains to prove the two nullity claims.  Fix an activation atom \(A_0\).
Formula (1) gives atomwise product Lebesgue measure on the cube.  For
distinct coordinates \(r,s\in V\), the hyperplane
\(\{\pi(r)=\pi(s)\}\cap[0,1]^V\) has product Lebesgue measure \(0\): by
Fubini, its section over any fixed value of all coordinates except \(r\) is
the singleton \(\{\pi(s)\}\), which has one-dimensional Lebesgue measure
\(0\).  Therefore (1) gives
\[
\nu\{A=A_0,\ 0\le\pi(q)\le1\hbox{ for all }q,\ \pi(r)=\pi(s)\}=0 .
\tag{3}
\]
There are only finitely many activation atoms.  Summing (3) over all atoms
with \(u,v\in A_0\) gives measure \(0\) for the displayed diagonal event
\(\pi(u)=\pi(v)\); the event in the statement requires only activated
coordinates to lie in the cube, and replacing that by the full cube changes
nothing by the atomwise full-cube support proved above.  Likewise, if
\(z\in U\cup T\), then \(z\ne u,v\).  Summing (3) with
\((r,s)=(z,u)\) over atoms containing \(u,v,z\) gives zero measure to
\(\pi(z)=\pi(u)\), and summing (3) with \((r,s)=(z,v)\) gives zero measure to
\(\pi(z)=\pi(v)\), on the same activation and cube support.  The boundary
event in the statement is contained in the union of these two null events, so
it also has \(\nu\)-measure \(0\).  This proves both remaining conclusions.
\end{proof}
